\documentclass[10.5pt,a4paper]{jsarticle}
\usepackage{graphicx}
\usepackage{enumitem}
% \usepackage{setspace}
% \setlength{\textwidth}{15mm} %% 本文の幅
% \setlength{\oddsidemargin}{} %% 左マージン
\setlength{\topmargin}{-0.4mm} %% 上部マージン
% \setlength{\btmmargin}{-0.4mm} %% 上部マージン
\setlength{\headheight}{0mm} %% 変更しない
\setlength{\headsep}{0mm} %% 変更しない
% \setlength{\textheight}{240mm} %% 本文の高さ(1)
\addtolength{\textheight}{\topskip} %% 本文の高さ(2)
\flushbottom
\pagestyle{empty} %% ページ番号つけない
%
%
\begin{document}

\begin{center}
\fontsize{12pt}{12pt}\bf{応用情報工学科 教室会議議事録}
\end{center}
\vspace{1\baselineskip}

%

\begin{itemize}[labelsep=0.5zw, leftmargin=1zw, itemindent=0pt, labelwidth=0pt, listparindent=0pt, parsep=0pt]
    \item[①] \textbf{日時：}2026年2月6日（金）　12:30 ～ 14:00
    \item[②] \textbf{場所：}応用情報工学科215会議室
    \item[③] \textbf{出席者（敬称略）：}%
    \parbox[t]{0.8\textwidth}{
        細野，香取，高橋，西脇，保谷，松野，望月，吉川，澤邉，滕，村上，山口，岡﨑，依田，関，岡崎，福田
    }
\end{itemize}

%全員 2013\item 出席者（敬称略）：作田，泉，木原，佐伯，三枝，作田，高橋(芳)，中川，中村，平山，山本，吉川，高野，岩田，塚本，高橋(聖)，細野，望月，芦澤，今池，五味，高橋(友)，井上，香取，山口
%全員 2017\item 出席者（敬称略）：吉川，泉，木原，高橋(聖)，平山，細野，香取，澤邉，松野，望月，五味，滕，関，高橋(友)，山口
%全員　2018\item 出席者（敬称略）：吉川，香取，木原，高橋(聖)，西脇，平山，細野，泉，澤邉，松野，望月，五味，滕，関，山口
%全員　2018\item 出席者（敬称略）：吉川，香取，木原，高橋(聖)，西脇，細野，泉，澤邉，松野，望月，五味，滕，関，山口
%全員　2019\item 出席者（敬称略）：吉川，香取，高橋(聖)，西脇，細野，泉，木原，澤邉，望月，五味，滕，村上，関，福田，山口
%全員　2026\item 出席者（敬称略）：細野，香取，高橋，西脇，保谷，松野，望月，吉川，五味，澤邉，滕，村上，山口，福田，岡﨑，依田，関

\vspace{0.1\baselineskip}

\begin{enumerate}[label=\arabic*：, leftmargin=2zw, labelsep=1zw, itemsep=1\baselineskip]
    \item \textbf{配布資料}
        % ・ の部分
        \begin{enumerate}[label=・\ 資料\arabic*：, leftmargin=6zw, nosep]
            \item 前々回議事録
            \item 前回議事録（案）
        \end{enumerate}

    \item \textbf{回覧資料}
        \begin{itemize}[label=・, leftmargin=1.5zw, nosep]
            \item 特になし．
        \end{itemize}
    
%%%%%%%%%%%%%%%%%%%%%%%%%%%%%%%%%%%%%%%%%%%%%%%%%%%%%%%%%%%%%%%%%%
%%%%%%%%%%%%%%%%%%%%%%%%%%%%%% 報告事項 %%%%%%%%%%%%%%%%%%%%%%%%%%%%%%%
%%%%%%%%%%%%%%%%%%%%%%%%%%%%%%%%%%%%%%%%%%%%%%%%%%%%%%%%%%%%%%%%%%
    
        \item \textbf{報告事項}
            
		\begin{itemize}[label=・, leftmargin=1.5zw, itemsep=1\baselineskip]%---------------
		\item 担当主任会議報告 (細野主任)
		\begin{itemize}[label=-, leftmargin=2zw, nosep]
			\item 
			%%%%%%%メモ
			教職員法律相談制度が設置された．
			学籍移動に関しては審議ではなく報告となる．ただし除籍は除く．休学は除く．
		\end{itemize}

		\item 各種募集 (細野主任)
		\begin{itemize}[label=-, leftmargin=2zw, nosep]
			\item 特になし．
		\end{itemize}

		\item 学生の状況
		\begin{itemize}[label=-, leftmargin=2zw, nosep]
			\item 1年生についてガイダンスは全員出席
			\item 後援会委員には習志野高校出身の学生の親御さんにお願いした
			\item 配慮あり学生がまだ登校できていないため，進路変更の相談をした
			\item 今年度から履修時期が変わり，取るべき科目を取れないなどの事象が発生した
			\item 転科の学生は一次抽選に参加できなかった
			\item フレッシュパーソンミーティングは無事終了した
			\item 2年生特になし
			\item 3年生について
			\item 4002秋葉　休学
			\item 昨年受賞したメタバースの学生が今年も受賞しオープンキャンパスで展示する
			\item 4年生について
			\item 3090皆川　休学または退学予定
			\item 井上　見つけ次第福田へ連れてくる
			\item 保護者面談の対応人員を登録してほしい→全学年対応済み
		\end{itemize}

		\item 入試・広報
		\begin{itemize}[label=-, leftmargin=2zw, nosep]
			\item 入試（五味准教授）
				\begin{itemize}
					\item 入試の状況について報告があった．
					\item 今年は偏差値を上げている．
				\end{itemize}
			\item 広報（松野教授）
				\begin{itemize}
				\item 学部のwebサイトが更新した．
				\item パンフレットの印刷代，webサイトの更新代が必要
			\end{itemize}
		\end{itemize}

		\item 就職（望月教授）
		\begin{itemize}[label=-, leftmargin=2zw, nosep]
			\item 2027年3月卒業の学生について，
			\item 報告するよう指導する．
			\item 院進学する予定の学生については松浦さんに教員が報告してほしい．
			\item 大学院について4名（20名中）のみ進路決定報告されている．
			\item 連休までが勝負なので，就活していない学生には指導する．
		\end{itemize}

		\item 各種委員会報告 
		\begin{itemize}[label=-, leftmargin=2zw, nosep]
			\item 委員会（山口教授）
			\begin{itemize}
				\item 予算
				\item 
			\end{itemize}
			
			\item 国際学術委員会（村上准教授）
			\begin{itemize}
				\item これから大学院生を増やすことを考えたときに海外の大学から進学されることを想定し，英語の案内へのURLをつけた
				\item 
			\end{itemize}

			\item 高大連携委員会（福田助教）
			\begin{itemize}
				\item 6名いるうち研究室を希望した3名を希望した研究室の先生へお願いした．残り3名は福田が見る．
				\item 
			\end{itemize}

		\end{itemize}

		\item その他
		\begin{itemize}[label=-, leftmargin=2zw, nosep]
			\item 五味准教授
			\begin{itemize}
				\item メールの添付ファイル非推奨の確認．
			\end{itemize}
			西脇教授
				今年度ゼミ生の募集をやめる

		\end{itemize}

	\end{itemize}

%%%%%%%%%%%%%%%%%%%%%%%%%%%%%%%%%%%%%%%%%%%%%%%%%%%%%%%%%%%%%%%%%%%%%%%
%%%%%%%%%%%%%%%%%%%%%%%%%%%%%% 審議事項 %%%%%%%%%%%%%%%%%%%%%%%%%%%%%%%
%%%%%%%%%%%%%%%%%%%%%%%%%%%%%%%%%%%%%%%%%%%%%%%%%%%%%%%%%%%%%%%%%%%%%%%

 \item \textbf{審議事項}
	\begin{itemize}[label=・, leftmargin=1.5zw, itemsep=1\baselineskip]%
		\item 前回議事録（案）について (細野主任)
		\begin{itemize}[label=-, leftmargin=2zw, nosep]
			\item 前回議事録（案）（資料2）について，指摘事項はなく承認された．
		\end{itemize}
		
		\item 学務（山口准教授）
		\begin{itemize}[label=-, leftmargin=2zw, nosep]
			\item 特になし．
			\item 松野先生がカリキュラム改定の資料を共有した
			\item 情報処理組み込みネットワークの三分野は残すべき
			\item AI関連の授業が少ない
			\item 組み込み系の授業が多い
			\item 学年間の連携を強化したい
			\item 話し合う時間を増やしていきたい
			\item データサイエンスはほぼすべての研究室で行っているので授業として増やしたい
			\item 卒業研究の懸け橋になるような授業を増やした方が良い
			\item 既存科目の名称変更，合併，設置時期の変更をした方が良い
			\item 既存科目名のまま内容の変更をしていきたい
			\item 三分野の名称は今の時代に合わせて変更してもいいのではないか
			\item 教室会議に合わせてカリキュラムを改定するためのワーキンググループを作りたい（松野）
			\item 方向性を決めてから検討していった方がいいのではないか（西脇）
			\item 学務経験者と実験担当の教員でワーキンググループを作れないか（松野）
			\item 実験項目のブラッシュアップのしていきたい（望月）
			\item カリキュラム改定に向けて話し合った方がよく，直近では実験間の差を埋めることを目標にしたい（松野）
			\item 一年時の授業に回路の授業が多すぎる問題はある（望月）
			\item 科目要覧にある樹形図を考え直した方がいい（高橋）
			\item カリキュラム改定は令和10年度なので令和9年度の夏までに提案をまとめる必要がある（山口）
			\item 大学としての制約もあるので話し合いをした方が良い（望月）
			\item 高校の情報授業が変わったためにPCの実機を触ったことが無いのではないか（澤部）
			\item ワーキンググループについては継続審議となった
		\end{itemize}
		
		\item 大学院（保谷教授）
		\begin{itemize}[label=-, leftmargin=2zw, nosep]
			\item 研究指導計画書を提出するよう指導してほしい（保谷）．
			\item TAの承認者が変更になった
			\item 履修登録確認表は学生に配らなくてよく，教員が確認しておく．

		\end{itemize}
			早期卒業なしで提出済み
			特待生もGPA最も高い学生を提出済み
			次年度以降は特待生の可能性がある学生が学部生以外であった場合の決定方法は継続審議とする
			学長賞も単純な成績だけで決定するかどうかは継続審議とする

			大学院の推薦基準はGPAの半分以上のものとして提出した．

			五味先生
			事務支援サイトにてs616も予約できるようにできないか
			→否決された
			朝日新聞のデジタル購読について申し込みが27名しかいないため，良いアイディアが欲しい
			→教員がまず登録し，研究室学生へ展開→その後学部生にアナウンスをすればいいのではないか
			資格等奨励寄付のアナウンスを今年もしていいか
			→承認
			学科で契約している
			→継続審議
			ワクチンソフトのインストール台数が131台あるが，131台契約してもよいか
			→五味先生に必要な台数を連絡し契約してもらう

			高橋先生
			専攻内説明か15時10分からでもいいか
			→承認された
		\item その他
		\begin{itemize}[label=-, leftmargin=2zw, nosep]
			\item -	委員会（教授）
			\begin{itemize}
				\item 
			\end{itemize}

		\end{itemize}
	\end{itemize}

	\item 次回の予定
	\begin{itemize}[label=・, leftmargin=1.5zw, nosep]
		\item 5月　21日（木）12：30〜
	\end{itemize}

\end{enumerate}

\begin{flushright}
以  上
\end{flushright}
\end{document}

%%%%%%%%%%%%%%以下メモ
